\section*{Capítulo \ref{ch:g6:decimals} --- Números decimales}

\begin{solution}{exo:g6:decimals:1}
$5.37$; \quad $0.9$; \quad $12.004$; \quad $8.05$.
\end{solution}

\begin{solution}{exo:g6:decimals:2}
Dígito de las décimas: $3$; dígito de las centésimas: $5$; el $8$ cuenta las decenas.
\[
86.354 = 8 \times 10 + 6 + 3 \times \frac{1}{10} + 5 \times
\frac{1}{100} + 4 \times \frac{1}{1000}.
\]
\end{solution}

\begin{solution}{exo:g6:decimals:3}
$5.3 > 5.29$ (comparar $5.30$ con $5.29$); \quad
$0.7 = 0.70$ (final \omterm{def:g1:counting:zero}{ceros} no cambiar nada); \quad
$2.09 < 2.9$ (décimas: $0 < 9$); \quad
$14.5 > 14.49$.
\end{solution}

\begin{solution}{exo:g6:decimals:4}
Rellenar con dos decimales: $4.20$; $4.05$; $4.51$; $4.15$; $4.50$. Orden:
\[
4.05 < 4.15 < 4.2 < 4.5 < 4.51 .
\]
\end{solution}

\begin{solution}{exo:g6:decimals:5}
Cada tick es una décima parte. $ A $: $6.3$; \quad $ B $: $6.7$; \quad
$ C $: $7.2$.
\end{solution}

\begin{solution}{exo:g6:decimals:6}
$45.80 + 7.65 = 53.45$.

$23.40 - 8.72 = 14.68$ (controlar: $14.68 + 8.72 = 23.4$).

$56 \times 24 = 1344$, y dos dígitos decimales en los factores:
$5.6 \times 2.4 = 13.44$.
\end{solution}

\begin{solution}{exo:g6:decimals:7}
$6.42 \times 10 = 64.2$; \quad
$0.35 \times 1000 = 350$; \quad
$78.1 \div 100 = 0.781$; \quad
$5 \div 1000 = 0.005$.
\end{solution}

\begin{solution}{exo:g6:decimals:8}
$2.4$ metro $= 240$ centímetro; \quad $370$ gramo $= 0.37$ kilogramos (dividir por $1000$);
\quad $0.85$ kilómetros $= 850$ metro; \quad $56$ milímetros $= 0.056$ metro.
\end{solution}

\begin{solution}{exo:g6:decimals:9}
A la unidad: $24$ (siguiente dígito $8$: \omterm{def:g4:numbers:round}{redondo} arriba). Al décimo: $23.9$.
Hasta el centésimo: $23.87$. Y $9.97$ hasta el décimo: el siguiente dígito es
$7$, entonces \omterm{def:g4:numbers:round}{redondo} hasta el $9$ décimas --- $10.0$.
\end{solution}

\begin{solution}{exo:g6:decimals:10}
Precio de las baguettes: $3 \times 1.15 = 3.45$.
Cambiar:
$5 - 3.45 = 1.55$.
\end{solution}

\begin{solution}{exo:g6:decimals:11}
Entre $7.4$ y $7.5$: por ejemplo $7.45$ (o $7.41$, $7.499$,
\dots).
Entre $3.99$ y $4$: por ejemplo $3.995$. En ambos casos
hay \emph{infinitamente muchos} tales números: siempre se pueden sumar más
lugares decimales.
\end{solution}

\begin{solution}{exo:g6:decimals:12}
Mayor: coloque los dígitos más grandes primero y el punto tan tarde como
posible: $85.2$. Más pequeño: los dígitos más pequeños primero y el punto como
lo antes posible: $2.58$.
\end{solution}

\begin{solution}{pb:g6:decimals:1}
\textbf{1.} $3$ es más grande: acolchado, $3.000 > 2.999$. El
\omterm{ex:g1:subtraction:difference}{diferencia} es $3.000 - 2.999 = 0.001$, una milésima.

\textbf{2.} estrictamente entre $7.4 = 7.40$ y $7.5 = 7.50$ mentir
\[
7.41,\ 7.42,\ 7.43,\ 7.44,\ 7.45,\ 7.46,\ 7.47,\ 7.48,\ 7.49 :
\]
nueve números, uno por cada nueva marca de graduación.

\textbf{3.} La misma imagen, diez veces más pequeña: entre $7.440$
y $7.450$ mentir $7.441, 7.442, \dots, 7.449$ --- nueve números
otra vez.

\textbf{4.} Entre $7.400$ y $7.500$, los números con tres
los dígitos decimales son $7.401, 7.402, \dots, 7.499$: todos los
milésimas de $401$ a $499$, eso es $99$ números.

\textbf{5.} La receta del zoom: escribe ambos números con el mismo
número de decimales, entonces \emph{añade un decimal más}
--- el número más pequeño seguido de un dígito de $1$ a $9$ tierras
strictly between the two. En efecto $5.67 = 5.670$ y
$5.68 = 5.680$, y
\[
5.670 < 5.675 < 5.680 ;
\]
asimismo $0.1999 = 0.19990 < 0.19995 < 0.20000 = 0.2$. entre
dos graduaciones vecinas siempre hay un nivel completamente nuevo de
nueve graduaciones más finas.

\textbf{6.} $3 < 3.05 < 3.1$; entonces $3 < 3.005 < 3.01$; entonces
$3 < 3.0005 < 3.001$. Cada candidato es superado por un zoom más.

\textbf{7.} Cualquiera que sea el número que proponga Tom, llámalo suyo.
candidato, estrictamente mayor que $3$ --- la pregunta 5 produce una
número estrictamente entre $3$ y el candidato. Entonces el candidato
era \emph{no} el número más cercano a $3$: algo \omterm{def:g3:numbers:evenodd}{incluso} más cerca
existe. Dado que esto le sucede a todos los candidatos sin excepción,
ningún número puede ser "el número justo después" $3$'': simplemente lo hace
no existe.

\textbf{8.} $2.9 < 2.95 < 3$, \ $2.99 < 2.995 < 3$, \
$2.999 < 2.9995 < 3$. Y $3 - 2.999 = 0.001$. con veinte
nueves, el \omterm{ex:g1:subtraction:difference}{diferencia} es un \omterm{def:g4:decimals:point}{coma decimal} seguido por diecinueve
\omterm{def:g1:counting:zero}{ceros} y un $1$ --- una unidad con el vigésimo decimal:
pequeño, pero no \omterm{def:g1:counting:zero}{cero}. Por muchos nueves que escriba Tom, nunca
alcanza $3$.

\textbf{9.} Los números enteros aumentan en pasos de $1$: más que $7$
pero menos que $8$ significaría un número entero de unidades estrictamente
entre $7$ y $8$ unidades, y no hay ninguna. El zoom se escapa
esto solo escribiendo dígitos \emph{después del punto} --- exactamente
lo que los números enteros no tienen.

\textbf{10.} La longitud más pequeña se redondea a $3.7$ es $3.65$ centímetros:
su siguiente dígito es $5$, cual \omterm{def:g4:numbers:round}{rondas} \emph{arriba}
(\cref{def:g6:decimals:rounding}). Y $3.75$ \omterm{def:g4:numbers:round}{rondas} arriba a $3.8$
por la misma razón. Entonces ``$3.7$ cm a la décima más cercana'' significa:
al menos $3.65$ cm, y estrictamente menos de $3.75$ cm --- un entero
Un segmento ampliado de posibles longitudes reales se esconde detrás de uno.
valor redondeado.

\textbf{11.} $0.0999 < 0.123 < 0.58 < 0.6$. Una frase:
comparando dígito por dígito desde la izquierda, $0.0999$ tiene $0$ décimas
mientras $0.6$ tiene $6$, y ahí se resuelve la comparación --- el
El número de dígitos escritos no dice nada sobre el tamaño.
(\cref{met:g6:decimals:compare}).

\textbf{12.} No: en centavos, $4.99$ euros es $499$ centavos y $5$
euros es $500$ centavos --- consecutivos \emph{entero} números, con
nada entre ellos. Los precios, con exactamente dos decimales,
son en realidad números enteros de centavos disfrazados: como el conjunto
números de la pregunta 9, sí tienen vecinos de al lado. el
El zoom sin fin necesita el derecho de escribir cada vez más decimales.

\textbf{13.} Los candidatos cerca $6$ son $5.82$ y $8.25$ (un
número que comienza con $2$, o con $58$, $25$, $82$, $85$, está lejos
de $6$). Distancias: $6 - 5.82 = 0.18$ y $8.25 - 6 = 2.25$.
Lo más cercano es $5.82$.

\textbf{14.} Por ejemplo $3.1415$, desde
$3.1410 < 3.1415 < 3.1420$; entonces $3.14159$, desde
$3.14150 < 3.14159 < 3.14160$. (Ambos son pasos reales en el
búsqueda real de $\pi = 3.14159\dots $)

\textbf{15.} Supongamos que alguien nombra un número tan grande como
como. Un zoom convierte cada par de graduaciones vecinas
entre $0$ y $1$ en nueve números nuevos y el zoom se puede
repetido para siempre --- las preguntas 2, 3 y 4 fueron solo las primeras
dos niveles. Repetirlo suficientes veces produce más \omterm{def:g6:decimals:places}{números decimales} entre $0$ y $1$ que el número nombrado. entonces no hay número
es lo suficientemente grande como para contarlos: hay infinitos.
\end{solution}
